\begin{definition}[不同制备事件的 PINEM 系综]
令 $f(\mathbf R,t_0)$ 是电子横向位置和到达时刻的归一化概率密度。
每次事件在自己的轨迹上经历耦合 $\beta(\mathbf R,t_0)$；未记录事件标签时，
输出为相应纯态的经典混合，能谱为
\begin{equation}
\overline P_\ell=\int f(\mathbf R,t_0)
J_\ell^2(2|\beta(\mathbf R,t_0)|)\,\dd^2R\dd t_0.
\label{eq:feqo-ensemble-pinem}
\end{equation}
\label{chap:feqo-averaging}
这不同于同一电子内部包络切片之间的相干叠加；后者必须先求总振幅。
\end{definition}

\begin{exbox}[束腰与光脉冲宽度怎样改变边带人口]
取二维高斯束 $f_R(R)=\exp[-R^2/(2\sigma_R^2)]/(2\pi\sigma_R^2)$，
场积分的横向剖面为 $\beta(R)=\beta_0\exp[-R^2/(2w_\beta^2)]$。
求弱驱动下的平均第一边带概率，再加入独立的高斯到达时刻抖动，比较结果与轴上准时入射。

\textbf{解。} 若 $|\beta_0|\ll1$，由 $J_1(2b)=b+O(b^3)$ 得
\begin{align*}
\overline P_1&=|\beta_0|^2\int_0^\infty
\frac{R\dd R}{\sigma_R^2}
\exp\left[-\left(\frac1{2\sigma_R^2}+\frac1{w_\beta^2}\right)R^2\right]
+O(|\beta_0|^4)\\
&=\frac{|\beta_0|^2}{1+2\sigma_R^2/w_\beta^2}+O(|\beta_0|^4).
\end{align*}
若还独立有高斯到达时刻分布，标准差为 $\sigma_t$，并取
$\beta(t_0)\propto\exp[-t_0^2/(2\tau_L^2)]$，则多出一维积分因子
$[1+2\sigma_t^2/\tau_L^2]^{-1/2}$。
例如 $\sigma_R=w_\beta$、$\sigma_t=\tau_L$ 时，总第一边带人口
是轴上准时电子的 $1/(3\sqrt3)$ 倍，而不是某个平均振幅的平方。
强场时应使用原 Bessel 积分，不能继续沿用这个低阶概率。
\end{exbox}

\begin{proposition}[强度平均与参考相位平均读取不同态分量]
若所有事件有相同 $b=|\beta|$，只有相位 $\phi$ 随事件改变，则
\[
\overline\rho_{\ell m}=J_\ell(2b)J_m(2b)
\langle\ee^{\ii(\ell-m)\phi}\rangle.
\]
所以能谱保持 $J_\ell^2(2b)$。圆周均匀相位使所有非对角元为零；
参考附近标准差为 $\sigma_\phi$ 的高斯相位使其乘
$\exp[-(\ell-m)^2\sigma_\phi^2/2]$。
\end{proposition}
\begin{proof}
逐事件的振幅为 $c_\ell=J_\ell(2b)\ee^{\ii\ell\phi}$。
对 $c_\ell c_m^*$ 作概率平均即得第一式；两个分布的 Fourier 特征函数给出后两式。
因此随机参考相位可以完全抹去后续聚束与相位读出，却不必改变单次能谱人口。
\end{proof}
